\documentclass[a4paper,11pt]{article}

\usepackage[utf8]{inputenc}
\usepackage[T1]{fontenc}
\usepackage[french]{babel}
\usepackage{geometry}
\usepackage{hyperref}
\usepackage{listings}
\usepackage{xcolor}

\geometry{margin=2.5cm}

\title{Tutoriel : Utiliser l'API Spotify avec Angular et Node.js}
\author{}
\date{\today}

\definecolor{codegray}{gray}{0.95}

\lstset{
    backgroundcolor=\color{codegray},
    basicstyle=\ttfamily\small,
    breaklines=true,
    frame=single
}

\begin{document}

\maketitle

\tableofcontents
\newpage

\section{Introduction}

Ce tutoriel explique comment utiliser l'API Spotify afin de recuperer des playlists a partir d'une recherche utilisateur dans le projet SkyBeat.

L'architecture utilisee est la suivante :

\begin{itemize}
\item Frontend : Angular
\item Backend : Node.js avec Express
\item API externe : Spotify Web API
\end{itemize}

Le backend sert d'intermediaire entre l'application Angular et l'API Spotify afin de gerer l'authentification, le cache et les requetes.

\section{Prerequis}

Avant de commencer, assurez-vous d'avoir installe :

\begin{itemize}
\item Node.js
\item Angular CLI
\item npm
\end{itemize}

Creer egalement une application sur le portail developpeur Spotify :

\begin{center}
\url{https://developer.spotify.com/dashboard}
\end{center}

Une fois l'application creee, recuperer :

\begin{itemize}
\item Client ID
\item Client Secret
\end{itemize}

Puis configurer les variables dans \texttt{backend/.env} (\texttt{SPOTIFY\_CLIENT\_ID}, \texttt{SPOTIFY\_CLIENT\_SECRET}, \texttt{SPOTIFY\_REDIRECT\_URI}).

\section{Architecture du projet}

Structure recommandee :

\begin{lstlisting}
SkyBeat/
    backend/
        serveur.js
        data/
            playlistCache.json
            users.json
    src/
        environments/
            environment.ts
        app/
            services/
                spotify.service.ts
\end{lstlisting}

\section{Backend Node.js}

Installer les dependances :

\begin{lstlisting}
cd backend
npm install
\end{lstlisting}

\subsection{Serveur Express}

Le fichier utilise dans SkyBeat est \texttt{backend/serveur.js}.

\begin{lstlisting}
app.get("/spotify/search", async (req, res) => {
  try {
    const { mood } = req.query;

    if (!mood) {
      return res.status(400).json({ error: "Missing mood parameter" });
    }

    const cacheKey = `spotify:${mood}`;
    const cached = getCachedOrNull(cacheKey);
    if (cached) {
      const normalizedCached = normalizeSpotifyPlaylists(cached);
      return res.json(normalizedCached);
    }

    const accessToken = await getSpotifyAccessToken();

    const response = await spotifyRequestWithRetry(
      () => axios.get("https://api.spotify.com/v1/search", {
        headers: { Authorization: `Bearer ${accessToken}` },
        params: { q: String(mood), type: "playlist", limit: "10" }
      }),
      "search_playlists"
    );

    const playlists = response.data.playlists.items
      .filter((item) => item && item.id)
      .map((item) => ({
        playlistId: item.id,
        title: item.name || "Untitled playlist",
        owner: item.owner?.display_name || "Spotify",
        trackCount: item.tracks?.total ?? 0,
        thumbnail: item.images[0]?.url || "",
        description: item.description || "",
        spotifyUrl: item.external_urls?.spotify || `https://open.spotify.com/playlist/${item.id}`,
      }));

    const normalizedPlaylists = normalizeSpotifyPlaylists(playlists);
    setCache(cacheKey, normalizedPlaylists);
    res.json(normalizedPlaylists);
  } catch (error) {
    res.status(500).json({ error: "Failed to search Spotify" });
  }
});
\end{lstlisting}

\subsection{Service Spotify}

La gestion du token Client Credentials est aussi faite dans \texttt{backend/serveur.js}.

\begin{lstlisting}
async function getSpotifyAccessToken() {
  if (spotifyAccessToken && Date.now() < spotifyTokenExpiry) {
    return spotifyAccessToken;
  }

  const clientId = process.env.SPOTIFY_CLIENT_ID;
  const clientSecret = process.env.SPOTIFY_CLIENT_SECRET;

  const authString = Buffer.from(`${clientId}:${clientSecret}`).toString("base64");

  const response = await spotifyRequestWithRetry(
    () => axios.post(
      "https://accounts.spotify.com/api/token",
      "grant_type=client_credentials",
      {
        headers: {
          Authorization: `Basic ${authString}`,
          "Content-Type": "application/x-www-form-urlencoded",
        },
      }
    ),
    "client_credentials_token"
  );

  spotifyAccessToken = response.data.access_token;
  spotifyTokenExpiry = Date.now() + (response.data.expires_in - 300) * 1000;

  return spotifyAccessToken;
}
\end{lstlisting}

\section{Frontend Angular}

Dans SkyBeat, le frontend Angular existe deja dans le dossier racine \texttt{src/}.

Lancer l'application :

\begin{lstlisting}
npm install
npm start
\end{lstlisting}

\subsection{Service Angular}

Le service utilise est \texttt{src/app/services/spotify.service.ts}.

\begin{lstlisting}
@Injectable({ providedIn: 'root' })
export class SpotifyService {
  private readonly http = inject(HttpClient);
  private readonly apiUrl = environment.backendUrl;

  searchPlaylists(mood: string): Observable<SpotifyPlaylist[]> {
    return this.http
      .get<SpotifyPlaylist[]>(`${this.apiUrl}/spotify/search`, {
        params: { mood },
      });
  }

  createPlaylist(title: string, trackIds: string[], oauthToken: string): Observable<SpotifyPlaylistResponse> {
    return this.http.post<SpotifyPlaylistResponse>(
      `${this.apiUrl}/spotify/playlist`,
      { title, trackIds },
      { headers: { Authorization: `Bearer ${oauthToken}` } }
    );
  }
}
\end{lstlisting}

\subsection{Composant de recherche}

Exemple d'appel depuis un composant Angular :

\begin{lstlisting}
import { Component } from '@angular/core';
import { SpotifyService } from '../services/spotify.service';

@Component({
  selector: 'app-generate',
  templateUrl: './generate.component.html'
})
export class GenerateComponent {
  mood = 'chill vibes music';
  playlists: any[] = [];

  constructor(private spotifyService: SpotifyService) {}

  search() {
    this.spotifyService.searchPlaylists(this.mood)
      .subscribe((data: any[]) => {
        this.playlists = data;
      });
  }
}
\end{lstlisting}

\subsection{Template HTML}

\begin{lstlisting}
<input [(ngModel)]="mood" placeholder="Mood Spotify (ex: chill, dance, jazz)">
<button (click)="search()">Search</button>

<ul>
  <li *ngFor="let playlist of playlists">
    {{ playlist.title }} - {{ playlist.owner }}
  </li>
</ul>
\end{lstlisting}

\section{Resultat}

L'utilisateur peut :

\begin{itemize}
\item saisir un mood musical
\item lancer une recherche via Angular
\item recuperer des playlists Spotify depuis le backend Express
\end{itemize}

\section{Ameliorations possibles}

\begin{itemize}
\item Ajouter la persistance des tokens utilisateur (Redis / session store)
\item Ajouter un refresh token flow complet pour l'OAuth Spotify
\item Ajouter une pagination cote frontend
\item Ajouter une vue detail d'une playlist (image, nombre de titres, description)
\end{itemize}

\end{document}
